\documentclass[10pt,twocolumn]{article}
\usepackage[T1]{fontenc}
\usepackage[utf8]{inputenc}
\usepackage{lmodern}
\usepackage{microtype}
\usepackage{amsmath}
\usepackage{amssymb}
\usepackage[margin=0.75in]{geometry}
\usepackage[hidelinks]{hyperref}
\usepackage{graphicx}
\usepackage{booktabs}
\usepackage{caption}
\usepackage{natbib}
\captionsetup{font=small,skip=3pt}
\setlength{\textfloatsep}{8pt plus 1pt minus 2pt}
\setlength{\floatsep}{8pt plus 1pt minus 2pt}
\setlength{\intextsep}{8pt plus 1pt minus 2pt}
\setlength{\abovedisplayskip}{5pt plus 1pt minus 2pt}
\setlength{\belowdisplayskip}{5pt plus 1pt minus 2pt}
\setlength{\abovedisplayshortskip}{4pt plus 1pt minus 1pt}
\setlength{\belowdisplayshortskip}{4pt plus 1pt minus 1pt}

\title{PAnDa: Parallel-block Adaptive Contrast DoLa}
\author{}
\date{}

\begin{document}
\maketitle

\begin{abstract}
Hallucination remains a central obstacle to deploying large language models in factual settings \citep{ji2023survey}. Inference-time decoding remains attractive because it can improve truthfulness without retrieval or retraining. DoLa showed that contrasting deeper and shallower layer signals can improve factuality with modest overhead \citep{chuang2023dola}, while later work suggests that factual risk is context-sensitive rather than uniform across a generation \citep{wan2023sequence,wastl2025token,zhang2025actlcd}. This note asks a narrower question: if fixed-strength layer contrast is already a strong baseline, can adaptivity be moved inside block-parallel decoding rather than implemented as full-response reranking?

We answer with \textbf{PAnDa}, a \emph{Parallel-block Adaptive Contrast DoLa} decoder that keeps two fixed contrast regimes inside each speculative block, detects the first local disagreement point between them, and applies truth-biased arbitration only from that point onward. The method is adaptive, but the adaptation is local and structured: it does not search over full completed responses, and it does not require a separate response-level confidence module. On the current TruthfulQA sanity-\(10\) development artifact, PAnDa is the strongest saved configuration in the project, improving over Greedy, DoLa, and a response-level alpha-switch reranker on the main TruthfulQA multiple-choice metrics. This development note presents initial findings and positions PAnDa as a promising direction for structured, adaptive factual decoding.
\end{abstract}

\section{Introduction}
Inference-time contrastive decoding is appealing because it can improve factuality without retrieval, tool use, or model retraining \citep{ji2023survey,lewis2020retrieval}. DoLa is the natural starting point: it contrasts deeper and shallower transformer layers and shows that layer-level internal disagreement can be useful at decoding time \citep{chuang2023dola}. Our own project then found that a simpler \textbf{DoLa-FixedAlpha} branch was a stronger practical baseline than heavier adaptive-\(\alpha\) variants. That result shifted the main question. The issue was no longer whether fixed contrast can work; it was whether adaptivity can be reintroduced in a way that stays narrow, interpretable, and computationally disciplined.

One response is full-response reranking: generate a conservative candidate, generate a more truth-seeking candidate, then choose between them. This is conceptually clean, but it is coarse. It treats truth-biased decoding as a search problem over completed answers, even when only a small local region is truly ambiguous. It also duplicates substantial decoding work, because multiple full candidates must be materialized before any adaptive decision can be made.

This note takes a different path. We propose \textbf{PAnDa}, a \emph{Parallel-block Adaptive Contrast DoLa} decoder that moves adaptive contrast selection inside block-parallel refinement. Instead of constructing multiple complete answers and comparing them after the fact, PAnDa refines a single evolving trajectory. At each speculative block, it computes a low-\(\alpha\) safe view and a high-\(\alpha\) truth-seeking view from the same shared hidden-state computation, detects where those two views first disagree, and arbitrates locally only from that point onward.

The resulting method should be read as a block-parallel adaptive extension of DoLa-FixedAlpha. Its claim is not that it globally optimizes the whole response. The claim is narrower: adaptive truth pressure can be made more local, more structured, and potentially more efficient when it is formulated as \emph{convergence-time arbitration inside a block} rather than as reranking among full candidates.

\section{Empirical Motivation: Contrast Strength}
\paragraph{Empirical effect of contrast strength.}
Before turning to background and method, it is useful to state the empirical pattern that motivated PAnDa in the first place. Table~\ref{tab:contrast-sweep} isolates the current saved baselines for Greedy decoding, standard DoLa, and three fixed-\(\alpha\) DoLa variants. These rows come from the saved fixed-\(\alpha\) sweep artifact on the same TruthfulQA sanity-\(10\) slice, so the cleanest comparison is in the quality and truthfulness metrics rather than raw wall-clock runtime. On this slice, all methods tie at \(MC1 = 0.10\), so the clearest changes appear in the softer truthfulness-oriented metrics \(MC2\) and \(MC3\).

\begin{table}[t]
\centering
\footnotesize
\begin{tabular}{lccc}
\toprule
Decoder & \(\alpha\) & MC2 & MC3 \\
\midrule
Greedy & -- & 0.1045 & 0.0500 \\
DoLa & adaptive & 0.1003 & 0.0250 \\
DoLa-FixedAlpha & 0.50 & 0.1992 & 0.1083 \\
DoLa-FixedAlpha (Low) & 0.10 & 0.1030 & 0.0500 \\
DoLa-FixedAlpha (High) & 0.95 & 0.2023 & 0.1417 \\
\bottomrule
\end{tabular}
\caption{Effect of contrast strength on the current TruthfulQA sanity-\(10\) development slice. Higher fixed contrast improves the soft truthful-ranking metrics \(MC2\) and \(MC3\), while the hard top-choice metric \(MC1\) remains flat across all rows.}
\label{tab:contrast-sweep}
\end{table}

The trend in Table~\ref{tab:contrast-sweep} is the key empirical motivation for PAnDa. Weak contrast behaves almost like no contrast at all. Greedy decoding reaches \(MC2 = 0.1045\), and the low-\(\alpha\) fixed variant stays nearly identical at \(MC2 = 0.1030\) with the same \(MC3 = 0.0500\). Standard DoLa does not help on this slice either, which suggests that merely invoking layer contrast is not enough if the effective contrast signal is too diffuse or too rigid. By contrast, the stronger fixed settings change the ranking behavior substantially: moving to \(\alpha = 0.50\) nearly doubles \(MC2\), and \(\alpha = 0.95\) pushes both \(MC2\) and \(MC3\) higher still.

This pattern suggests a precise interpretation of contrast and truthfulness. Stronger contrast does not improve the hard winner-take-all metric \(MC1\) on this small slice, so it should not be described as improving answer accuracy in a broad absolute sense. What it does improve is \emph{truth preference}: more probability mass shifts toward true options, and more true options outrank the strongest false distractor. In other words, increasing contrast primarily improves \emph{truthful ranking quality} before it improves top-1 answer selection. That is exactly the regime where local arbitration becomes attractive. If high contrast is often useful for truthfulness but too aggressive to apply everywhere, the natural next step is to keep both a conservative and a truth-seeking regime available and switch locally only when the block state provides evidence for doing so.

\section{Background: From Fixed Contrast to Local Arbitration}
The project logic behind PAnDa is simple. DoLa demonstrates that layer contrast can improve factuality \citep{chuang2023dola}. Follow-up work such as SLED and Speculative Contrastive Decoding broadens the contrastive-decoding design space \citep{zhang2024sled,yuan2023scd}. At the same time, sequence-level certainty and token-level self-consistency both suggest that hallucination risk is uneven across a generation \citep{wan2023sequence,wastl2025token}. Active layer-contrastive control sharpens that lesson further: a single always-on contrast regime is often too rigid \citep{zhang2025actlcd}.

The most direct adaptive fix is response-level candidate selection, but that framing has two weaknesses. First, it is late: by the time the selector acts, both full responses have already been generated. Second, it is coarse: the method can keep candidate \(A\) or candidate \(B\), but it cannot preserve the good prefix of one and correct only the ambiguous suffix of the other. PAnDa is motivated by the idea that this is the wrong granularity for many factual questions. If only a small portion of a continuation needs stronger truth pressure, then adaptation should happen there rather than at the level of whole completed answers.

This motivates a convergence-based interpretation. Instead of asking, ``which final answer is better?'', PAnDa asks, ``where inside the current speculative block do safe and truth-seeking contrast regimes first meaningfully diverge, and how should that divergence alter the block update?'' That question is local enough to exploit block-parallelism, but still adaptive enough to avoid a single rigid contrast setting everywhere.

\section{Method: A Block-Parallel Arbitration Pipeline}
\begin{figure*}[t]
\centering
\fbox{%
\parbox{0.96\textwidth}{%
\small
\textbf{PAnDa pipeline.}
Prompt and current prefix
\(\rightarrow\)
initialize speculative block \(Y^{(0)}\)
\(\rightarrow\)
run a shared forward pass over the block
\(\rightarrow\)
build safe and truth-seeking contrastive views
\(\rightarrow\)
measure local disagreement
\(\rightarrow\)
find the earliest divergence point
\(\rightarrow\)
apply truth-biased local arbitration from that point onward
\(\rightarrow\)
commit the stable prefix
\(\rightarrow\)
repeat until the continuation is complete.
}}
\caption{PAnDa as a local arbitration pipeline. The decoder maintains one evolving trajectory, computes safe and truth-seeking contrastive views from a shared block forward pass, detects the first meaningful divergence point, and applies truth-biased arbitration only from that point onward.}
\label{fig:panda-pipeline}
\end{figure*}

Figure~\ref{fig:panda-pipeline} summarizes the method.

\paragraph{Stage 1: Shared block forward pass.}
Let the next \(W\) tokens be represented by a speculative block
\begin{equation}
Y^{(k)} = [y_t^{(k)}, y_{t+1}^{(k)}, \ldots, y_{t+W-1}^{(k)}],
\end{equation}
where \(k\) indexes the Jacobi refinement pass. Instead of advancing one token at a time, PAnDa runs a shared forward pass over the current block estimate and extracts logits for all positions in that block at once. The method still uses Jacobi-style refinement internally: like a classical Jacobi iteration, it forms the next block estimate from the previous block iterate in parallel across positions, then repeats refinement until a stable prefix can be committed.

\paragraph{Stage 2: Dual-regime contrastive views.}
For each block position, PAnDa computes two fixed-strength contrastive views from the same final and shallow logits:
\begin{equation}
\begin{aligned}
z_{\text{safe}} &= z_{\text{final}} - \alpha_{\text{low}} z_{\text{shallow}}, \\
z_{\text{truth}} &= z_{\text{final}} - \alpha_{\text{high}} z_{\text{shallow}}.
\end{aligned}
\end{equation}
In the current Stage 12 artifact, \(\alpha_{\text{low}}=0.1\) and \(\alpha_{\text{high}}=0.95\). The intent is not to estimate an optimal \(\alpha\) continuously. The intent is to expose a conservative and a corrective regime from the same block state.

\paragraph{Stage 3: Local disagreement detection.}
For each position \(i \in \{1,\ldots,W\}\), PAnDa measures whether the two regimes disagree materially. In the current implementation, a position counts as a true local disagreement only when the low-\(\alpha\) and high-\(\alpha\) top-1 tokens differ and the Jensen-Shannon divergence between their local predictive distributions exceeds a threshold. Concretely,
\begin{equation}
d_i = \mathrm{JSD}\!\left(
\mathrm{softmax}\!\left(z_{\text{safe}}[i] / \tau\right),
\mathrm{softmax}\!\left(z_{\text{truth}}[i] / \tau\right)
\right).
\end{equation}
If no position with a token mismatch exceeds the divergence threshold, the block update is accepted normally. If some positions do exceed it, PAnDa finds the earliest meaningful divergence point,
\begin{equation}
i^{*} = \min \{ i : d_i > \tau_{\text{div}} \},
\end{equation}
and activates arbitration only from that point onward.

\paragraph{Stage 4: Truth-biased local arbitration.}
From the first divergence point onward, PAnDa decides whether the truth-seeking local update is sufficiently better than the safe local update. In abstract form,
\begin{equation}
\hat{z}_i =
\begin{cases}
z_{\text{truth}}[i], & \text{if } s_{\text{truth}}(i) > s_{\text{safe}}(i) - \beta, \\
z_{\text{safe}}[i], & \text{otherwise,}
\end{cases}
\end{equation}
where \(\beta\) is a truth-bias margin. This is the key methodological move. PAnDa does not choose between completed candidates; it chooses between local updates inside the same evolving block.

\paragraph{Stage 5: Prefix commit and refinement.}
After arbitration, the block is refined again. If the block stabilizes, PAnDa commits the stable prefix and continues. If not, it still commits the maximum stable prefix it can justify and then repeats the refinement process on the suffix. This makes PAnDa neither plain autoregressive decoding nor full-response search. It is block-local approximate search with selective truth-biased intervention.

\section{Experiments and Development Results}
\paragraph{Benchmarks and scope.}
Our long-term project benchmark set remains TruthfulQA and HaluEval. This note focuses on the clearest current development artifact: the Stage 12 TruthfulQA sanity-\(10\) run. The purpose of this slice is diagnostic rather than final. It is small enough to inspect manually, but large enough to expose the characteristic behavior of the decoder.

\paragraph{Baselines.}
We compare PAnDa against Greedy decoding, standard DoLa, and a response-level alpha-switch reranker built from low-\(\alpha\) and high-\(\alpha\) fixed-contrast candidates. All runs use the same self-calibrated 1.5B backbone and the same TruthfulQA multiple-choice sampling slice.

\begin{table}[t]
\centering
\footnotesize
\begin{tabular}{lccc}
\toprule
Decoder & MC2 & MC3 & Latency (s) \\
\midrule
Greedy & 0.1045 & 0.0500 & 4.26 \\
DoLa & 0.1003 & 0.0250 & 7.14 \\
Alpha-Switch Reranker & 0.2023 & 0.1417 & 14.32 \\
PAnDa & \textbf{0.2669} & \textbf{0.1667} & 17.77 \\
\bottomrule
\end{tabular}
\caption{TruthfulQA sanity-\(10\) development results from the current Stage 12 artifact. All methods tie at \(MC1 = 0.10\), so the main differences appear in \(MC2\), \(MC3\), and latency.}
\label{tab:panda-main}
\end{table}

Table~\ref{tab:panda-main} summarizes the current development outcome. PAnDa is the strongest saved configuration on this slice. The most important point is not that it dominates every metric by a huge margin; the most important point is that the block-local arbitration story appears viable. Relative to standard DoLa, the method has a much stronger truthful-ranking profile. Relative to the response-level alpha-switch reranker, it improves the two main TruthfulQA multiple-choice ranking metrics, raising \(MC2\) from \(0.2023\) to \(0.2669\) (about a \(32\%\) relative increase) and \(MC3\) from \(0.1417\) to \(0.1667\).

The interpretation is consistent with the earlier contrast-strength sweep. Plain DoLa is too rigid: it applies one contrastive behavior everywhere. A response-level alpha-switch reranker is adaptive, but it is late and coarse, because it can only keep one full candidate. PAnDa occupies the middle ground. It can preserve the shared good prefix, localize where stronger truth pressure becomes relevant, and only apply the corrective regime from that point onward. That is the core reason its current artifact is compelling.

\section{Discussion}
\paragraph{Why this framing is attractive.}
The most important advantage of PAnDa is not just that it is adaptive, but that it is adaptive at the right granularity. Many factual questions do not require a fully different response; they require a local correction, a safer reformulation, or a stronger refusal to commit to a plausible falsehood in only a small part of the answer. Response-level reranking is too coarse for that pattern. PAnDa addresses it directly by making adaptation conditional on local divergence inside a speculative block.

\paragraph{Why this is still not a global optimizer.}
PAnDa should not be misunderstood as a full-response global search method. It does not enumerate all possible answers, nor does it globally optimize a full-sequence objective. It remains a heuristic decoding algorithm. Its claim is narrower and, we think, more defensible: \emph{for block-local ambiguity, local arbitration is a better computational match than reranking completed responses.}

\paragraph{Efficiency caveat.}
The current implementation still pays real refinement cost. On this slice, PAnDa is slower than the response-level reranker (\(17.77\)s versus \(14.32\)s), so the current paper should not make a strong speed claim. The right interpretation is a quality-latency trade-off: the additional refinement cost is buying measurably better truthful ranking quality, especially on \(MC2\) and \(MC3\). That is encouraging because it suggests the extra block passes are not wasted, but it also makes the next optimization target clear. Future work should focus on reducing refinement overhead by improving convergence and commit efficiency, so that more of the truthfulness gain can be retained at lower runtime.

\paragraph{What remains to validate.}
The present note is still a development document. The current Stage 12 artifact is the strongest saved configuration in the repo, but the evidence still comes from a small development slice. The next validation step is straightforward: keep the current settings fixed, scale the TruthfulQA evaluation, add a held-out protocol where possible, and confirm that the present quality-latency balance survives beyond the sanity-\(10\) run.

\section{Conclusion}
This note presents PAnDa, a block-parallel adaptive DoLa decoder that performs truth-biased local arbitration during Jacobi refinement rather than after full-response generation. The method keeps the strong fixed-contrast backbone, exposes both safe and truth-seeking local views from the same shared block state, and intervenes only from the first meaningful divergence point onward.

The broader contribution is conceptual. Adaptive contrastive decoding does not need to be framed as token-level control or as search over completed responses. A middle ground exists: local convergence-time arbitration. In the current project state, Stage 12 PAnDa is the strongest saved artifact and therefore the right method to carry forward. The remaining task is no longer to invent a new adaptive mechanism; it is to validate whether the present block-local mechanism remains the best choice on larger and more rigorous evaluations.

\bibliographystyle{plainnat}
\bibliography{references}

\end{document}
